\section{Related Work}
\label{sec:related_work}

\subsection{Digital Twins in Autonomous Laboratories}
The integration of robotic automation has significantly accelerated experimental throughput in materials science and chemistry. Pioneering works, such as mobile robotic chemists \cite{burger2020mobile} and self-driving thin-film laboratories \cite{macleod2020self}, have demonstrated the potential of automated execution. Recent autonomous laboratory platforms such as A-Lab \cite{szymanski2022automated} have further demonstrated end-to-end materials discovery. However, these systems predominantly rely on hard-coded waypoints. To mitigate the prohibitive costs of deploying new workflows in physical bio-laboratories, Digital Twin (DT) technologies have emerged as crucial architectural foundations \cite{dt_migration_mobile}. Complex bio-assays require high-fidelity contact physics rather than mere kinematic bounding-box approximations. To address this, our framework adopts the GPU-accelerated NVIDIA Isaac Sim \cite{makoviychuk2021isaac}, providing a mathematically rigorous tensor-based prior for simulating contact-rich operations.

\subsection{Safe Navigation and Fluid-Spill Prevention}
In safety-critical laboratory settings handling volatile reagents or biological cultures, ensuring strictly collision-free and spill-proof mobile base navigation is paramount. Standard sampling-based planners, such as RRT* or A*, excel at finding obstacle-free paths but often neglect task-space orientation constraints, rendering them unsuitable for transporting open liquid containers \cite{gasparetto2015path}. Recent work on digital twin-enabled path planning for mobile robots \cite{dt_port_inspection, dt_unity_ros} has shown promising results in structured environments. To guarantee absolute fluid stability, we formulate the macroscopic inter-station transfer as a constrained classical optimization problem. By enforcing strict roll-pitch bounds on the end-effector and minimizing the trajectory jerk integral, our classical planning module systematically eliminates the risk of liquid spillage prior to any physical deployment.

\subsection{Deep Reinforcement Learning for Robotic Manipulation}
Deep Reinforcement Learning (DRL) has shown tremendous promise in mastering contact-rich manipulation tasks. Residual RL approaches \cite{johannink2019residual} combine model-based controllers with learned residuals, while large-scale data collection pipelines \cite{levine2018learning} have demonstrated the feasibility of learning hand-eye coordination. Dexterous manipulation with demonstrations \cite{rajcsy2018learning} and in-hand manipulation \cite{andrychowicz2020learning} further showcase the capability of DRL for complex tasks. However, training robust DRL policies in reality is constrained by sample inefficiency and hardware wear-and-tear. Our framework utilizes massively parallel simulation to overcome this bottleneck, training PPO agents across thousands of simultaneous virtual environments with posture-aware reward shaping \cite{ng1999policy}.

\subsection{Sim-to-Real Transfer and Domain Randomization}
Bridging the sim-to-real gap remains a central challenge in robotic learning. Domain randomization \cite{tobin2017domain} has proven effective for visual robustness, while CAD-to-RL pipelines \cite{sadeghi2017cad} enable zero-shot transfer from simulation. Recent surveys \cite{hofer2021sim2real, zhao2020sim2real} have systematically categorized sim-to-real approaches across manipulation, locomotion, and navigation. For contact-rich tasks, system identification and adaptive dynamics randomization \cite{hwangbo2019learning} have shown significant improvements. Our approach combines extensive domain randomization with regularized few-shot real-world fine-tuning, achieving robust deployment with minimal physical demonstrations.

\subsection{Imitation Learning for Manipulation}
Imitation Learning (IL) provides a complementary paradigm to DRL by leveraging expert demonstrations. Behavioral Cloning (BC) and Generative Adversarial Imitation Learning (GAIL) \cite{ho2016gail} represent two mainstream approaches \cite{osuga2018survey}. More recently, Diffusion Policy \cite{chi2023diffusion} has demonstrated superior performance in visuomotor tasks through action denoising, while vision-language-action models \cite{brohan2023rt2} have shown the potential of transferring web knowledge to robotic control. Our framework adopts a BC-based pipeline with regularized few-shot fine-tuning, which is particularly well-suited for precision micro-manipulations where demonstration quality is critical.

\subsection{Hierarchical Task Execution and Skill Composition}
Long-horizon manipulation motivates decomposition into reusable skills. Hierarchical reinforcement learning (HRL) can learn high-level skill selection \cite{pateria2021hierarchical}, while rule-based state machines provide explicit transition guards and failure handling. Digital-twin-based task replanning has also been investigated \cite{dt_task_replanning}. Our current system uses the latter, deterministic approach: a hierarchical state machine sequences independently developed skills. We do not claim a trained HRL policy or a probabilistic meta-controller.
